\section{Receding-horizon model-predictive control}\label{sec:mpc}

This section is the conceptual core of the document. It develops the
finite-horizon stochastic optimal control problem, the receding-horizon
principle, the various stochastic and output-feedback variants, and then
specialises everything to the FSA-v2 closed-loop model-predictive
controller in both mathematical and implementation form. A closing
subsection sketches a linear--quadratic--Gaussian approximation that
could be added as a cheap baseline.

\subsection{Finite-horizon stochastic optimal control}\label{sec:mpc-finite}

Consider a discrete-time stochastic system on a state $X_k \in \Xcal$ and
a control $U_k \in \Ucal$, with transition kernel
$X_{k+1} \sim p(\,\cdot\mid X_k,\,U_k)$, a stage cost
$\ell:\Xcal \times \Ucal \to \R$, and a terminal cost
$V_f:\Xcal \to \R$. The \emph{open-loop} finite-horizon problem of length
$H$ is
\begin{equation}
  \min_{u_{0:H-1}\in\Ucal^H}
    \;\E\!\left[\,\sum_{k=0}^{H-1} \ell(X_k,u_k)
                  + V_f(X_H)\,\Big|\,X_0 = x\right],
  \label{eq:open-loop-fh}
\end{equation}
with the expectation taken over the joint law of $X_{1:H}$ given the
chosen sequence $u_{0:H-1}$. Denote its optimal value by $V_H(x)$ and a
minimising sequence by $u_{0:H-1}^*(x)$. The mapping
$x \mapsto u_0^*(x)$ is the open-loop optimal first-action policy.

For FSA-v2 the state $X_k = (B_k,F_k,A_k)$ lives in
$\Xcal = [0,1]\times[0,\infty)^2$, and the ``control sequence''
is the eight-dimensional RBF-coefficient vector $\thetactrl$ that
parameterises the entire schedule $\Phi(\,\cdot\,;\thetactrl)$ over
$[0,T]$ via \eqref{eq:schedule-decoder}. The discrete-time abstract
notation \eqref{eq:open-loop-fh} therefore maps onto the continuous
schedule of Section~\ref{sec:control} by a one-to-one identification
between admissible $u_{0:H-1}$ and admissible $\thetactrl$. The
expectation in \eqref{eq:open-loop-fh} is the same one already written
out in \eqref{eq:cost}.

\subsection{The receding-horizon principle}\label{sec:mpc-receding}

In \emph{receding-horizon} or \emph{model-predictive} control, the
open-loop problem \eqref{eq:open-loop-fh} is solved repeatedly on a
sliding window. At each replan time $t_n$, given a current state
estimate $\xhat_n$, one
\begin{enumerate}
  \item solves \eqref{eq:open-loop-fh} for $u_{0:H-1}^*(\xhat_n)$;
  \item applies only the first action $u_0^*(\xhat_n)$ over a stride of
        length $\Delta < H$;
  \item observes new data, updates the state estimate to $\xhat_{n+1}$,
        and returns to step~1.
\end{enumerate}
The resulting closed-loop policy is a function of $\xhat_n$ alone, so the
controller is implementable. Stability of the closed loop is delicate.
Mayne and Rawlings (2000) show that under suitable assumptions on the
horizon length and the terminal cost / terminal set, the closed-loop
state remains in a positively invariant set and converges to a
neighbourhood of a desired equilibrium. For FSA-v2 we rely on the
\emph{long-horizon} variant of that theory: the planning horizon
$H$ is chosen long enough relative to the slowest time constant
$\tau_B = 42$ days that the truncation error from neglecting trajectory
beyond $H$ is small.

\subsection{Stochastic and output-feedback variants}\label{sec:mpc-variants}

Three flavours of stochastic MPC are commonly distinguished.

\begin{enumerate}
  \item[(i)] \emph{Certainty-equivalent} MPC. Replace the expectation in
        \eqref{eq:open-loop-fh} by evaluation at a point estimate
        $\hat X = \E[X\mid \xhat_n]$, and solve the resulting deterministic
        problem. Cheap, but discards uncertainty information.
  \item[(ii)] \emph{Expected-cost} (stochastic) MPC. Keep the
        expectation, and approximate it by Monte Carlo (sampling
        trajectories of the stochastic dynamics) or by scenario trees.
  \item[(iii)] \emph{Output-feedback} MPC. When $X_k$ is not observed
        directly, the controller acts on a state estimate $\xhat_n$
        produced by an external estimator (Kalman, particle filter,
        SMC${}^2$, etc.). For nonlinear non-Gaussian systems no
        \emph{separation principle} holds in general --- the joint
        design of the estimator and the controller is part of the
        modelling, not a free composition.
\end{enumerate}

The FSA-v2 pipeline is in class (ii) and (iii) simultaneously. The state
is unobserved (only the four-channel proxies of
\eqref{eq:obs-hr}--\eqref{eq:obs-steps} are available), so a state
estimator is required: this is the rolling-window SMC${}^2$ filter of
Section~\ref{sec:smc2}. The expectation in the cost \eqref{eq:cost} is
kept, evaluated by Monte Carlo on the SDE under common random numbers,
and minimised in $\thetactrl$-space by the control SMC${}^2$ of
Section~\ref{sec:control}. The full closed-loop policy is therefore the
composition
\[
  \text{(observations)}
    \xrightarrow{\text{rolling SMC${}^2$ filter}} \xhat_n
    \xrightarrow{\text{control SMC${}^2$, posterior mean}} \bar\theta
    \xrightarrow{\eqref{eq:schedule-decoder}} \Phi(\,\cdot\,;\bar\theta)
    \xrightarrow{\text{plant}} \text{(new observations)}.
\]
There is no separation principle here. The fact that this composition
works at all is an empirical claim, validated in
Section~\ref{sec:results}.

\subsection{The FSA-v2 MPC pipeline --- mathematical statement}\label{sec:mpc-pipeline}

We now write down the closed-loop algorithm. Let $\Delta_w = 96$ bins be
the filter window length (one day), $\Delta_s = 48$ bins the stride
(twelve hours), and $H_T = 14 \cdot 96 = 1344$ bins the macrocycle
(fourteen days), so that the number of windows is
$N_w = (H_T - \Delta_w)/\Delta_s + 1 = 27$.

\begin{algorithm}[H]
\caption{FSA-v2 closed-loop MPC}
\label{alg:fsa-mpc}
\begin{algorithmic}[1]
\State Initialise the plant at $(B_0,F_0,A_0)$; initialise an empty
       observation buffer $Y_{1:0} = \emptyset$.
\For{$n = 1,\dots,N_w$}
  \State \textbf{Filter update.} If $n=1$, run cold-start SMC${}^2$
         (Algorithm~\ref{alg:smc2}) on $Y_{(n-1)\Delta_s+1:(n-1)\Delta_s+\Delta_w}$
         to produce a posterior over $\thetadyn$ for the first window;
         otherwise warm-start from the previous window's posterior via
         the Schr\"odinger--F\"ollmer bridge of Section~\ref{sec:smc2}.
  \State \textbf{State extraction.} For a small set of $\thetadyn$
         particles drawn from the posterior, run the bootstrap filter
         (Algorithm~\ref{alg:sir}) forward to the stride endpoint and
         compute the importance-weighted mean of the resulting
         $(B,F,A)$; call this $\xhat_n$.
  \State \textbf{Control plan.} Run the control SMC${}^2$ of
         Section~\ref{sec:control}, conditioned on $\xhat_n$ and the
         posterior-mean dynamics parameters $\bar\theta_{\mathrm{dyn}}$,
         to obtain the posterior-mean schedule coefficient
         $\bar\theta_{\mathrm{ctrl}}$.
  \State \textbf{Apply.} Decode
         $\Phi(\,\cdot\,;\bar\theta_{\mathrm{ctrl}})$ via
         \eqref{eq:schedule-decoder} and apply only its first stride
         (length $\Delta_s$) to the plant; collect the resulting
         observations and append them to the buffer.
\EndFor
\end{algorithmic}
\end{algorithm}

The horizon over which the controller plans (line~5) is the full macrocycle
remainder $H_T - (n-1)\Delta_s$, but only the first $\Delta_s$ bins of the
plan are committed (line~6). This is the receding-horizon truncation of
Section~\ref{sec:mpc-receding} applied to the FSA-v2 problem. Replanning
occurs every twelve hours, at the wake-window boundary so that the
morning-loaded burst envelope of \eqref{eq:phi-envelope} is delivered in
full each day.

\subsection{Code-level structure of the MPC implementation}\label{sec:mpc-code}

This subsection departs from the usual mathematical idiom and describes
how the abstract loop of Section~\ref{sec:mpc-pipeline} is realised in
the source code of the project. The reader who only wants the
mathematics may skip to Section~\ref{sec:mpc-feasibility}; the reader
who intends to run, modify, or extend the controller will find this
useful. We name three artefacts.

\paragraph{(a) The plant interface, \texttt{StepwisePlant}.} The
``plant'' in stochastic optimal control is whatever object the
controller acts upon to produce new observations. In the code base it is
the class \texttt{StepwisePlant} in
\href{run:../version_2/models/fsa_high_res/_plant.py}{\texttt{version\_2/models/fsa\_high\_res/\_plant.py}}.
A \texttt{StepwisePlant} carries the hidden true state $(B,F,A)$, the
true dynamics parameters, and a global bin counter; its single mutating
method is
\begin{quote}\small
\verb|advance(stride_bins, Phi_daily)|
$\;\to\;$ \texttt{(trajectory, obs\_HR, obs\_sleep, obs\_stress, obs\_steps)}.
\end{quote}
A call to \texttt{advance} expands the supplied daily-$\Phi$ array into
the morning-loaded sub-daily envelope of \eqref{eq:phi-envelope}
(preserving $\int_{\text{day}}\Phi\,\dd t = 24\,\Phi_d$ exactly), runs
Euler--Maruyama for \texttt{stride\_bins} fifteen-minute steps, samples
the four-channel observations of \eqref{eq:obs-hr}--\eqref{eq:obs-steps},
and updates the internal state. Pedagogically, the
\texttt{StepwisePlant} is the digital embodiment of the ``real world''
against which the controller is evaluated; mathematically, a single call
samples one transition of the joint state-observation process under the
chosen control over a stride of length $\Delta_s$.

\paragraph{(b) The MPC driver,
\texttt{bench\_smc\_full\_mpc\_fsa.py}.} The receding-horizon
orchestrator lives in
\href{run:../version_2/tools/bench_smc_full_mpc_fsa.py}{\texttt{version\_2/tools/bench\_smc\_full\_mpc\_fsa.py}}.
Its execution sequence implements Algorithm~\ref{alg:fsa-mpc} step by
step. For each window~$n$ it
\begin{enumerate}
  \item runs the rolling-window filter by calling
        \verb|run_smc_window| (cold-start, $n=1$) or
        \verb|run_smc_window_bridge| with a Schr\"odinger--F\"ollmer
        base measure (warm-start, $n\ge 2$);
  \item extracts the posterior-mean dynamics parameters and the
        smoothed latent state $\xhat_n$ at the end of the window, the
        latter computed by running Algorithm~\ref{alg:sir} forward on a
        small set of $\thetadyn$ particles drawn from the posterior and
        taking the weighted mean of the resulting $(B,F,A)$ samples;
  \item assembles a \texttt{ControlSpec} (cost functional, horizon
        length, current state estimate, RBF basis);
  \item calls the control SMC${}^2$ loop in
        \href{run:../smc2fc/control/tempered_smc_loop.py}{\texttt{smc2fc/control/tempered\_smc\_loop.py}},
        which returns a posterior-mean schedule coefficient
        $\bar\theta_{\mathrm{ctrl}}$ and a corresponding $\Phi$
        schedule;
  \item calls \verb|plant.advance(stride_bins, Phi_daily)| to apply
        only the first stride of that schedule;
  \item appends the resulting observations to the data buffer and
        shifts the window.
\end{enumerate}

\paragraph{(c) The posterior-mean handoff between windows.} The smoothed
state $\xhat_n$ extracted in step~2 above is what initialises the next
window's particle filter. Consecutive windows are therefore connected
\emph{at the state level} by a single point estimate (the importance-weighted
smoothed mean) and \emph{at the parameter level} by the SF
Bures--Wasserstein bridge of Section~\ref{sec:smc2}. This is the
practical realisation of the warm-start that gives the rolling SMC${}^2$
its computational tractability: at fourteen days of fifteen-minute data
the cold-start cost would be $27\times$ larger.

\medskip
\noindent
Table~\ref{tab:mpc-files} collects the mapping between the steps of
Algorithm~\ref{alg:fsa-mpc} and the source files that realise them.

\begin{table}[h]
\centering
\small
\begin{tabular}{@{}lll@{}}
\toprule
Step in Algorithm~\ref{alg:fsa-mpc} & Concept & Source \\
\midrule
3 & rolling SMC${}^2$ filter      & \texttt{smc2fc/core/tempered\_smc.py} \\
3 & SF Bures--Wasserstein bridge   & \texttt{smc2fc/core/sf\_bridge.py} \\
3 & inner bootstrap filter         & \texttt{smc2fc/filtering/gk\_dpf\_v3\_lite.py} \\
4 & smoothed-state extraction      & \texttt{version\_2/tools/bench\_smc\_full\_mpc\_fsa.py} \\
5 & control SMC${}^2$ outer loop   & \texttt{smc2fc/control/tempered\_smc\_loop.py} \\
5 & cost functional, RBF decoder   & \texttt{version\_1/models/fsa\_high\_res/control.py} \\
6 & plant advance                   & \texttt{version\_2/models/fsa\_high\_res/\_plant.py} \\
2,6 & sub-daily $\Phi$ envelope    & \texttt{version\_2/models/fsa\_high\_res/\_phi\_burst.py} \\
\bottomrule
\end{tabular}
\caption{Mapping from Algorithm~\ref{alg:fsa-mpc} to the source code.}
\label{tab:mpc-files}
\end{table}

\subsection{Recursive feasibility and the F-barrier}\label{sec:mpc-feasibility}

The fatigue ceiling $\Fmax = 0.40$ is enforced through the soft-barrier
penalty $\lambda_F \cdot \max(F-\Fmax,0)^2$ in the cost
\eqref{eq:cost}, not through a hard constraint. The choice is
deliberate. Under stochastic dynamics with state-dependent diffusion and
unbounded driving noise, an exact hard constraint $F(t) \le \Fmax$ cannot
be guaranteed by any controller without making the admissible-control
set empty for some achievable initial conditions; this is the standard
\emph{recursive feasibility} difficulty in stochastic MPC. A soft
penalty preserves recursive feasibility -- there is always a feasible
schedule -- and produces the constraint approximately, with the violation
rate decreasing as $\lambda_F$ is increased. Empirically (see
Section~\ref{sec:results}) the closed-loop F-violation rate is
$1.28\%$ on the $84$-day open-loop benchmark and $0.00\%$ on the
$14$-day closed-loop benchmark, well within the $5\%$ acceptance gate.

\subsection{Robustness and the flat-cost-surface argument}\label{sec:mpc-robustness}

The controller of Section~\ref{sec:mpc-pipeline} commits the posterior
mean $\bar\theta_{\mathrm{ctrl}}$ rather than a sample from the control
posterior or the full distribution. One might worry that this discards
useful information, especially when the filter posterior over
$\thetadyn$ is uncertain or biased. Empirically the policy is
nevertheless robust to filter parameter drift, and the explanation is
geometric: near the optimum, the cost $J(\thetactrl)$ is locally flat in
the directions explored by the SMC${}^2$ cloud. A first-order Taylor
expansion of $J$ at the posterior mean gives, for a perturbed estimate
$\bar\theta + \delta$,
\begin{equation}
  J(\bar\theta + \delta)
   \;\approx\; J(\bar\theta) + \tfrac{1}{2}\delta^{\!\top} H \delta,
   \qquad H = \nabla^2 J(\bar\theta),
\end{equation}
so the additional cost from using $\bar\theta$ in place of an unknown
``oracle'' optimum scales with the Hessian. When $H$ is small in the
directions of greatest filter uncertainty, the closed-loop performance is
correspondingly insensitive to filter accuracy. Section~\ref{sec:results}
reports closed-loop performance on a record where only $3$ out of $26$
estimable parameters pass a strict coverage gate; the controller
nevertheless achieves a $+17\%$ amplitude gain. The flat-cost-surface
argument is the mechanism.

\subsection{An exact-approximate alternative: LQG and the differential
Riccati equation}\label{sec:mpc-lqg}

We close Section~\ref{sec:mpc} by sketching a second controller that
could be added to the project as a cheap baseline or as a warm-start for
the SMC${}^2$ control of Section~\ref{sec:control}. The sketch is
deliberately complete enough to make the assumptions and limitations
explicit; it is not yet implemented in the code base, and is flagged as
required next work in Section~\ref{sec:discussion-lqg}.

\paragraph{Linearisation around an operating point.}
Three assumptions reduce the FSA-v2 problem to one with a closed-form
optimal feedback law.

\begin{assumption}\label{ass:lqg-A1}
There is a nominal trajectory $(x^*(t),\Phi^*(t))$ on $[0,T]$, taken in
the simplest case to be the deterministic-skeleton equilibrium
$(B^*,F^*,A^*)(\Phi_{\mathrm{const}})$ under a constant control. With
$\tilde x = x - x^*$ and $\tilde \Phi = \Phi - \Phi^*$, define the
linearised drift coefficients
\begin{equation}
  A_{\mathrm{lin}}(t) = \frac{\partial f}{\partial x}\bigg|_{*},
  \qquad B_{\mathrm{lin}}(t) = \frac{\partial f}{\partial \Phi}\bigg|_{*},
  \label{eq:lqg-AB}
\end{equation}
where $f$ is the drift of \eqref{eq:fsa-B}--\eqref{eq:fsa-A}.
\end{assumption}

The cubic Stuart--Landau term $-\eta A^3$ is linearised away under
Assumption~\ref{ass:lqg-A1}, so the bistable / limit-cycle behaviour of
the autonomic amplitude is lost. The linearisation is valid only in a
neighbourhood of the operating point.

\begin{assumption}\label{ass:lqg-A2}
The state-dependent Jacobi and CIR diffusions in
\eqref{eq:fsa-B}--\eqref{eq:fsa-A} are replaced by constant-coefficient
Gaussian noise $\bar\sigma\,\dd W$ evaluated at the operating point.
\end{assumption}

The $[0,1]$ and $[0,\infty)$ state constraints are no longer enforced
implicitly by the noise vanishing at the boundary, so trajectories may
leave the physiological domain. Acceptable when fluctuations around the
operating point are small relative to the boundary distances.

\begin{assumption}\label{ass:lqg-A3}
The cost \eqref{eq:cost} is replaced by the quadratic surrogate
\begin{equation}
  J_{\mathrm{LQ}}(\tilde \Phi)
    \;=\; \E\!\left[\,\int_0^T \bigl(\tilde x^{\!\top} Q \tilde x
                       + \tilde \Phi^{\!\top} R \tilde \Phi\bigr)\,\dd t
                       + \tilde x(T)^{\!\top} Q_T \tilde x(T)\,\right],
  \label{eq:lqg-cost}
\end{equation}
with $Q \succeq 0$, $R \succ 0$, $Q_T \succeq 0$.
\end{assumption}

To approximate the $-\!\int A\,\dd t$ reward, one chooses
$x_{\mathrm{ref}} \neq x^*$ shifted toward higher $A$ and penalises
$\tilde x = x - x_{\mathrm{ref}}$ with $Q$ giving most weight to the $A$
component; this preserves $Q \succeq 0$. The soft fatigue barrier
$\lambda_F \cdot \max(F-\Fmax,0)^2$ is replaced by a quadratic in
$F-F^*$, which penalises both directions symmetrically -- too restrictive
on the rest side and too lax on the over-fatigue side near the barrier.

\paragraph{The differential Riccati equation.}
Under Assumptions~\ref{ass:lqg-A1}--\ref{ass:lqg-A3}, the optimal
feedback law for the linearised system is the time-varying linear policy
\begin{equation}
  \tilde \Phi^*(t) \;=\; -K(t)\,\tilde x(t),
  \qquad K(t) \;=\; R^{-1}\,B_{\mathrm{lin}}(t)^{\!\top}\,P(t),
  \label{eq:lqg-feedback}
\end{equation}
where $P:[0,T] \to \R^{3\times 3}$, $P(t) \succeq 0$, solves the
\emph{differential Riccati equation}
\begin{equation}
  -\dot P(t) \;=\; A_{\mathrm{lin}}(t)^{\!\top} P(t)
                   + P(t)\,A_{\mathrm{lin}}(t)
                   - P(t)\,B_{\mathrm{lin}}(t)\,R^{-1}\,
                     B_{\mathrm{lin}}(t)^{\!\top} P(t)
                   + Q,
  \qquad P(T) = Q_T.
  \label{eq:riccati}
\end{equation}
For a time-invariant linearisation (constant $A_{\mathrm{lin}}$,
$B_{\mathrm{lin}}$) and an infinite horizon, $P$ solves the
\emph{algebraic} Riccati equation
$0 = A_{\mathrm{lin}}^{\!\top} P + P\,A_{\mathrm{lin}} - P\,B_{\mathrm{lin}}\,R^{-1}\,B_{\mathrm{lin}}^{\!\top} P + Q$.
Under Assumptions~\ref{ass:lqg-A1}--\ref{ass:lqg-A2} the state-estimation
problem also linearises and is solved by the Kalman--Bucy filter; the
LQG controller is then the certainty-equivalent composition of
Kalman--Bucy filter and LQR feedback, and the (linear-Gaussian)
separation principle holds.

\paragraph{Where this would fit in the project.}
A complete LQG companion to the SMC${}^2$ controller would consist of a
new module \texttt{smc2fc/control/lqg/} containing
\begin{enumerate}
  \item automatic linearisation of the drift in
        \texttt{version\_2/models/fsa\_high\_res/\_dynamics.py} around a
        user-supplied operating point, via JAX automatic differentiation;
  \item a numerical solver for \eqref{eq:riccati} by backward integration
        on the same fifteen-minute time grid used elsewhere in the
        project;
  \item a \texttt{LQGController} class emitting the feedback gain $K(t)$
        of \eqref{eq:lqg-feedback} and a corresponding $\Phi$ schedule
        on the project's time grid.
\end{enumerate}
The natural uses are (i) as a \emph{cheap baseline} against which the
SMC${}^2$ controller of Section~\ref{sec:control} is benchmarked --- does
the non-linear posterior-mean controller actually outperform the linear
LQG approximation, and by how much? --- and (ii) as a \emph{warm-start}
for the control posterior, by least-squares projection of the LQG
schedule onto the RBF basis to give an initial $\thetactrl$ around which
the SMC${}^2$ then explores. None of this currently exists in the
repository; we return to it in Section~\ref{sec:discussion-lqg}.
